\section{Related Works}

EMG-based gesture classification has been extensively studied in the biomedical signal processing
literature. \citet{Phinyomark2012} conducted a comprehensive evaluation of time-domain EMG
features, showing that Root Mean Square (RMS) and Zero Crossing Rate (ZCR) are among the most
discriminative and computationally efficient features for gesture recognition. Their work provided
a systematic foundation for multi-domain feature extraction pipelines and directly motivates the
inclusion of time-domain features in this study alongside statistical, frequency-domain, and
entropy features. Reported classification accuracies for binary gesture tasks using time-domain
features ranged from 70--85\%, providing a direct reference point for evaluating the present
results.

Support Vector Machines with kernel functions have been consistently identified as strong
performers for EMG classification. \citet{Oskoei2007} benchmarked multiple classifiers on
upper-limb EMG data and found that SVM with an RBF kernel outperformed k-Nearest Neighbour and
Linear Discriminant Analysis in most conditions, particularly when class boundaries were
non-linear in feature space. Their reported SVM accuracies of 80--95\% were obtained with
multi-channel (6--10 channel) recordings, which explains the expected performance gap relative
to the single-channel setup used here. This result directly motivates the inclusion of SVM-RBF
as a primary classifier in this study.

Ensemble methods have demonstrated competitive performance in biosignal tasks. \citet{Atzori2016}
evaluated Random Forest on the NinaPro EMG database alongside deep learning models, finding it
competitive while offering greater interpretability. XGBoost, introduced by \citet{Chen2016}, is
a regularised gradient-boosted tree ensemble that achieves state-of-the-art results across
diverse benchmarks. Including both RF and XGBoost alongside SVM-RBF provides a representative
comparison of tree-based ensembles versus kernel-based approaches.

For feature selection, Kruskal-Wallis hypothesis testing has been applied in biomedical signal
processing to identify statistically discriminative features without assuming normally distributed
data --- an important consideration for entropy features \citep{Subasi2013}. Combined with PCA
for dimensionality reduction, this two-stage selection pipeline reduces the risk of overfitting
while preserving the most informative feature components.
